\documentclass[10pt,a4paper]{amsart}
\usepackage[margin=19mm]{geometry}
\usepackage{amsmath,amssymb,mathtools}
\usepackage[scr=rsfs,cal=euler]{mathalfa}
\usepackage{xfrac}
\usepackage[unicode,hidelinks,bookmarks=false]{hyperref}
\newcommand{\ie}{i.e.\ }
\newcommand{\cf}{cf.\ }
\newcommand{\resp}{respectively\ }
\newcommand{\Subsection}{\subsection}
\providecommand{\nobreakdash}{\nobreak\hbox{-}}
\setlength{\emergencystretch}{2em}
\multlinegap0pt
\usepackage{bm}
\usepackage{bbm}
\usepackage{dsfont}
\usepackage{xspace}
\usepackage{cancel}
\usepackage{mathtools}
\usepackage[shortlabels]{enumitem}
\usepackage{amsthm}
\makeatletter
\@ifundefined{defn}{\newtheorem{defn}{Defn}}{}
\@ifundefined{prop}{\newtheorem{prop}[defn]{Proposition}}{}
\@ifundefined{thm}{\newtheorem{thm}[defn]{Theorem}}{}
\makeatother
\title[The List-distinguishing chromatic number of graphs containin]{The List-distinguishing chromatic number of graphs containing only small complete bigraphs}
\author{Amitayu Banerjee}
\date{Source: arXiv:2509.11992v1 (CC BY 4.0)}
\begin{document}
\maketitle

\noindent\emph{Source and licence.} The List-distinguishing chromatic number of graphs containing only small complete bigraphs. Version arXiv:2509.11992v1 (CC BY 4.0), distributed under the Creative Commons Attribution 4.0 International licence, as recorded in the arXiv licence metadata for this exact version and shown on the abstract page https://arxiv.org/abs/2509.11992v1. Full rights notice: licenses/arxiv-2509.11992-CC-BY-4.0.txt.\par\smallskip

\noindent\textit{Editorial scope.} Counted unit: the Prop whose source label is \texttt{Proposition 4.4} in the article (source0.tex lines 696--703, source label \texttt{Proposition 4.4}), delivered together with the article's own complete original proof of it (source0.tex lines 705--777, one \texttt{proof} environment of 73 lines). No mathematics is written, repaired or re-derived: statement and proof are the article's own text line for line. Required context retained uncounted: Theorem 1.7 (lines 151--158); Theorem 4.1 (lines 630--646). Every definition, display and label of those lines is retained unchanged. Omitted from this package and not redistributed: 1--150, 159--629, 647--695, 704--704, 778--1139. Nothing inside a retained excerpt is omitted or altered: every retained body is delivered exactly as the article's lines are written, so no omission receipt of any kind accompanies this package. Citations: the retained excerpts carry no citation command at all, so no bibliography entry of the article is transcribed and no reference is redistributed. Reference adaptation: none was needed and none was made; every reference inside the delivered unit resolves inside it, so no unresolved reference or citation is produced. Printed slips kept literal: no printed word, symbol, hypothesis or number of the counted statement or of its proof was altered. Layout and class: the article's own class is \texttt{amsart} with packages including hyperref, enumerate, geometry, amsmath, amssymb, amsthm, amsfonts, graphicx, tikz; the wrapper is the lane's pinned \texttt{amsart} editorial wrapper at 10pt on A4, which restates the theorem-like environments the retained text needs and adds the article's own macro definitions that the retained text needs (none), with extra wrapper packages (bm, bbm, dsfont, xspace, cancel, mathtools, enumitem). The delivered numbering is the unit's own. Assets: no retained span contains a figure environment, a graphics-inclusion command, a PDF-page inclusion, a table or a TikZ picture; no image file is copied into this package, the package declares no asset dependency and no image file is redistributed with it. Licence: version arXiv:2509.11992v1 of this article is distributed under the Creative Commons Attribution 4.0 International licence, as recorded in the arXiv licence metadata for this exact version and shown on the abstract page https://arxiv.org/abs/2509.11992v1. A full rights notice is delivered at licenses/arxiv-2509.11992-CC-BY-4.0.txt. Characters: every retained excerpt is pure ASCII, so the delivered engine, XeLaTeX with its default Latin Modern text font, reports no missing character.
\begin{thm}\label{Theorem 1.7}{\em Fix $4\leq r\leq \frac{\Delta(G)}{2}$. Let $G$ be a $(n-2r+1)$-connected graph of order $n>2r$ such that
$G$ is a ($K_{r,r+1}, K_{r+1,r+1}$)-free graph.
Then, 
\begin{center}
$\chi_{D_{L}}(G)\leq 2\Delta(G)-  (3\lfloor\frac{(\Delta(G)+2)}{r+1}\rfloor-4)$.
\end{center}
}
\end{thm}

\begin{thm}\label{Theorem 4.1}
{\em Let $X$ be any graph on $a=n-2r+1$ vertices satisfying the following properties:
\begin{enumerate}[(i)]
    \item $\alpha(X)\le r-1$, 
    \item $\kappa(X)\ge a-b$, and
    \item $\Delta(X)\ge r+2$.
\end{enumerate}

Let $H$ be any graph on $b$ vertices with $\alpha(H)\le r-1$. Then $G=X\vee H$ satisfies the following:
\begin{enumerate}
    \item $G$ is $(n-2r+1)$-connected,
    \item $G$ contains no induced subgraph isomorphic to $K_{r,r+1}$ or $K_{r+1,r+1}$,
    \item $\Delta(G)\ge 3r+1$. Thus, $\displaystyle \frac{\Delta(G) + 2}{r + 1} \ge 3$.
\end{enumerate}
Consequently, by Theorem \ref{Theorem 1.7}, we obtain $\chi_{D_{L}}(G)\leq 2\Delta(G)-5$ which is a better bound for $\chi_{D_{L}}(G)$ than $\chi_{D_{L}}(G)\leq 2\Delta(G)-1$.
}
\end{thm}

\begin{prop}\label{Proposition 4.4}
{\em Suppose $a=2m\geq r+8$ is an even integer and $1\leq k\leq 5$ is any integer. 
%So $a=2m$ for some $m\in\mathbb{N}$. 
Let $D_{2m}=\langle \rho,\sigma : \rho^m=\sigma^2=1,\ \sigma\rho\sigma=\rho^{-1}\rangle$ 
be the dihedral group of order $2m$, $R=\{\rho^j:0\le j\le m-1\}$, $R\sigma=\{\rho^j\sigma:0\le j\le m-1\}$, and $F\subseteq R\sigma$ satisfy $|F|=m-k$.
If $S \;=\; (R\setminus\{e\}) \;\cup\; F$ is a generating set,
then $X=\mathrm{Cay}(D_{2m},S)$ satisfies $(i)-(iii)$ of Theorem \ref{Theorem 4.1}.}
\end{prop}

\begin{proof}
%We only show when $k=5$. The rest of the cases are similar.
Let $d=|S|=(m-1)+(m-k)=2m-(k+1)$. Then $X$ is a $d$-regular graph.

\begin{enumerate}[(i)]
    \item We write $D_{2m}=R\cup R\sigma$ as the disjoint union of cosets $R$ and $R\sigma$.
    Now, each of $R$ and $R\sigma$ induces a clique in $X$.
Take two distinct rotations \(\rho^i,\rho^j\in R\) with \(i< j\). Then
$
(\rho^i)^{-1}\rho^j=\rho^{j-i}\in R\setminus\{e\}\subseteq S,
$
so \(\rho^i\) and \(\rho^j\) are adjacent in $X$.
Similarly, take two distinct reflections \(\rho^i\sigma,\rho^j\sigma\in R\sigma\) with \(i> j\). 
We note that $\sigma^{2}=\{e\}$, so $\sigma^{-1}=\sigma$ and $\sigma\rho^{k}\sigma=\rho^{-k}$ for all integers $k\geq 1$.
Then
\[
(\rho^i\sigma)^{-1}(\rho^j\sigma)
= (\sigma^{-1}\rho^{-i})(\rho^j\sigma) 
= \sigma\rho^{-i}\rho^j\sigma
= \sigma\rho^{\,j-i}\sigma
= \rho^{-(j-i)}=\rho^{\,i-j}\in S.
\]

So, \(\rho^i\sigma\) and \(\rho^j\sigma\) are adjacent in $X$.
Thus, any independent set in $X$ can contain at most one rotation and at most one reflection. Hence $\alpha(X)\leq 2< r-1$ since $r> 6$.
    
\item We show $\kappa(X)\geq 2m-(2k+1)$. By Menger's theorem, it suffices to show that for any two distinct vertices
$u,v\in V(X)$, there are $2m-(2k+1)$ internally vertex-disjoint $u$--$v$ paths.
By the arguments in (i), $R$ and $R\sigma$ induces a clique $K_m$ in $X$. Moreover, we can see that each rotation has exactly $m-k$ neighbours in $R\sigma$ and each reflection has exactly $m-k$ neighbours in $R$.
Pick $u,v\in V(X)$.

\medskip\noindent\textbf{Case A: $u,v\in R$ are distinct rotations.}  
Write $u=\rho^i$, $v=\rho^j$, $i\ne j$. Since $R$ induces a clique $K_{m}$, there is an edge $uv$ joining $u$ and $v$. Let $P_{0}=\{uv\}$. 
Let $R\setminus\{u,v\}=\{\rho^{r_1},\dots,\rho^{r_{m-2}}\}$. For each $\rho^{r_t}\in R\backslash\{u,v\}$ consider the path
\begin{center}
$P_t= \{u\rho^{r_t}, \rho^{r_t}v\}$.    
\end{center}

We construct $m-2k$ disjoint paths each using a distinct reflection as internal vertex.
We note that the set of reflections
adjacent to $u$ is a translate $\rho^i F$, and those adjacent to $v$ is $\rho^j F$ for some $i$ and $j$. Since $|F|=m-k$, we have $\vert N_{X}(u)\cap N_{X}(v)\cap R\sigma\vert =\vert\rho^i F\cap \rho^j F\vert\geq m-2k$.\footnote{$N_{G}(x)$ denotes the open-neighborhood of $x$ in $G$.}
Therefore, there are at least $m-2k$ reflections adjacent to both $u$ and $v$.
For each such reflection $s\in \rho^i F\cap\rho^j F$, consider the path
\[
Q_s= \{us, sv\}.
\]

Then $\{P_{0}\}\cup\{P_{t}:\rho^{r_{t}}\in R\backslash\{u,v\}\}\cup \{Q_{s}:s\in \rho^i F\cap\rho^j F, N_{X}(u)=\rho^{i}F, N_{X}(v)=\rho^{j}F\}$ is a collection of $(m-1)+(m-2k)=2m-(2k+1)$ internally vertex-disjoint paths.
 
\medskip\noindent\textbf{Case B: 
$u,v\in R\sigma$ are distinct reflections.} 
This is analogous to Case A. One only needs to swap the roles of rotations and reflections.

\medskip\noindent\textbf{Case C: $u$ is a rotation and $v$ is a reflection (or vice versa).} \\
Without loss of generality, assume $u=\rho^i$ and $v=\rho^j\sigma$ for some $i$ and $j$. 

\begin{itemize}
  \item[] \textsc{Subcase C(A):} Suppose $u$ and $v$ are adjacent. Consider the edge $uv$ as a path.   
  Since $v$ has $m-k$ neighbors in $R$, there are $m-(k+1)$ neighbors of $v$ in $R\backslash \{u\}$.
  For each such rotation $\rho^{r}\in R\setminus\{u\}$ (that is a neighbor of $v$) take the path $P_{\rho^{r}}=\{u\rho^{r},\rho^{r}v\}$.
  Similarly, there are $m-(k+1)$ different reflections distinct from $v$ that are neighbors of $u$. For each such reflection $s$, take the path $P_{s}=\{us, sv\}$. 
  Altogether, these give $1+2(m-(k+1))=2m-(2k+1)$ internally vertex disjoint paths from $u$ to $v$.
  \vspace{2mm}

  \item[] \textsc{Subcase C(B):} If $u$ is not a neighbor of $v$, then similar to \textsc{Subcase C(A)}, we can obtain at least $2m-(2k+1)$ internally vertex disjoint paths from $u$ to $v$.
\end{itemize}

\medskip
By Menger's theorem, $\kappa(X)\ge 2m-(2k+1)=a-(2k+1)> a-b$ since $b=2r-1>2k+1$.
    
    \item Since $X$ is $d$-regular and $a\geq r+8$, we have $\Delta(X)=2m-(k+1)=a-(k+1)\geq r+2$.
\end{enumerate}
\end{proof}

\end{document}
